Power quality disturbances (PQDs) can affect the reliability of electrical power systems and must be detected accurately. This work investigates the use of machine learning for classifying nine PQD classes from eight features extracted using the S-transform and statistical analysis. First, class-specific feature subsets are identified through physical reasoning, graphical and statistical analysis, and binary classification experiments. The results show that different disturbances depend on different feature combinations, with features related to harmonic content, voltage amplitude, and high-frequency components being particularly informative. Next, decision tree, random forest, and multilayer perceptron classifiers are optimized using stratified five-fold cross-validation and evaluated on a held-out test set. A soft-voting ensemble combines the predictions of the three models. The ensemble achieves a macro F1-score of 0.9844 and a normal-class precision of 0.9412, outperforming the individual models on both metrics. The results demonstrate that the extracted features provide a strong basis for automated PQD classification and that combining complementary classifiers improves performance. However, the remaining confusion between normal signals, transients, and flicker, together with the normal-class precision, indicates that further validation and improvement are required before deployment in a real-world power system.
